% !TEX root = ../../tree_notes.tex
% =====================================================================
%  Chapter 5 review -- the 15-section revision contract (chapter_end.txt).
%  Revision-only: tables, one-liners, pointers back via \thmref/\probref/
%  \secref/\coderef/\propref/\defref/\exref/\figref. No derivations.
%  Summit: LC Hard 2322 Minimum Score After Removals on a Tree
%  (root it, take tin/tout, read every edge-deletion as an interval split).
% =====================================================================

\chaptersection{Chapter 5 Summary}

Four sections, one physical question: \emph{lay the branching tree out along a line so a
sub-question becomes an array-question}. The timestamps (5.1) decide \emph{when} each node
lives; the flattening (5.2) turns a whole subtree into one contiguous interval; the path and
depth views (5.3) read the same walk as a stack or a level array; the $2n{-}1$ tour (5.4)
turns a distance into a range minimum. Every entry below is a pointer, not a re-derivation.

\begin{center}
\renewcommand{\arraystretch}{1.25}
\begin{tabular}{@{}p{0.25\linewidth} p{0.47\linewidth} p{0.20\linewidth}@{}}
\toprule
\textbf{Result} & \textbf{Statement (with conditions)} & \textbf{Where} \\
\midrule
Entry/exit timestamps & One DFS stamps $\tin(v)$ on first-in and $\tout(v)$ on last-out; the
  pair is a half-open interval $[\tin(v),\tout(v))$. & \defref{c24-times} \\
Ancestor test in $O(1)$ & $u$ is an ancestor of $v$ iff $[\tin(v),\tout(v))$ fits inside
  $[\tin(u),\tout(u))$ -- two integer comparisons. & \thmref{c24-anc} \\
Bracket characterisation & Erasing labels from the DFS leaves a balanced bracket string; every
  balanced string rebuilds exactly one ordered tree. & \propref{c24-bracket} \\
Subtree $=$ range & Order nodes by $\tin$; then $\subtree(v)$ is the block
  $[\tin(v),\tin(v){+}\size(v))$ of the flat array. & \thmref{c24-range} \\
Static subtree sums & A prefix array over the flat array answers any subtree sum in $O(1)$;
  invertible aggregates only. & \propref{c24-prefix} \\
Root-path by in/out sweep & Push an aggregate on entry and undo on exit: the running value
  \emph{is} the root-to-node path sum; a difference sweep converts subtree-add to path-read.
  & \propref{c24-sweep} \\
Distance via LCA & $\dist(u,v)=\depth(u)+\depth(v)-2\depth{\lca(u,v)}$ on the unique $u$--$v$
  path. & \thmref{c24-dist} \\
LCA is an RMQ (preview) & $\lca(u,v)$ is the shallowest node of the $2n{-}1$ tour between the
  two first-occurrences -- proven in the LCA chapter. & \propref{c24-rmqpreview} \\
\bottomrule
\end{tabular}
\end{center}

\medskip
{\small\textbf{Algorithms from this chapter -- one DFS, four payoffs.}}
\begin{center}
\small
\renewcommand{\arraystretch}{1.2}
\setlength{\tabcolsep}{3pt}
\begin{tabular}{@{}p{0.19\linewidth} l l p{0.29\linewidth} p{0.155\linewidth}@{}}
\toprule
\textbf{Algorithm} & \textbf{Time} & \textbf{Memory} & \textbf{When to use (deciding phrase)} & \textbf{Code} \\
\midrule
Stamp $\tin/\tout$, test ancestry
  & $O(n){+}O(1)$ & $O(n)$
  & ``is $u$ above $v$?'' / order nodes by the walk
  & \coderef{c24-stamp} \\
Flatten $+$ subtree sum by prefix
  & $O(n)$ pre, $O(1)$ q
  & $O(n)$
  & many ``sum/count over the whole subtree of $v$''
  & \coderef{c24-flatten-build} \\
Root-path stack / level array
  & $O(n)$ & $O(n){+}O(h)$
  & aggregate along a root-to-node path or per depth
  & \coderef{c24-path-stack} \\
$2n{-}1$ tour $+$ distance
  & $O(n)$ build & $O(n)$
  & ``distance / directions between two named nodes''
  & \coderef{c24-tour-build} \\
Summit: subtree XOR $+$ interval split
  & $O(n^2)$ & $O(n)$
  & ``cut two edges, score the three pieces''
  & \coderef{c24-summit} \\
\bottomrule
\end{tabular}
\end{center}

\chaptersection{Master Decision Flowchart}

Post-study ladder: walk top to bottom, stop at the first branch that matches. It decides
\emph{which reading of the tour} to build -- a follower never stamps the wrong clock.

\begin{lstlisting}[language={}, basicstyle=\ttfamily\small, numbers=none, frame=single]
  The tree is rooted. What does the QUESTION reach for?
  |
  |-- "is u an ancestor of v?" / "list nodes in DFS order" / nested-bracket strings
  |     +-- stamp BOTH-tick tin/tout; O(1) containment test            [5.1]
  |     +-- bracket <-> tree: push on '(', pop on ')'                  [5.1, LC 1600]
  |
  |-- "for each v, SUM/COUNT/MAX over its whole subtree" (many queries)
  |     +-- entry-only clock tin in [0,n), tout = tin + size
  |         subtree(v) = flat range [tin, tin+size)                    [5.2]
  |     +-- static  -> prefix array (invertible only)                  [5.2]
  |     +-- updates -> hand the range to Fenwick [c14] / segtree [c13]
  |
  |-- "sum along the path FROM THE ROOT to v" / "an answer per depth level"
  |     +-- carry a path value down, UNDO on the way up                [5.3]
  |     +-- bucket nodes by depth -> an ordinary array, then scan/sort
  |         (LC 1161 heaviest level, LC 2583 kth-largest level sum)    [5.3]
  |
  |-- "distance / walk between TWO named nodes"
  |     +-- dist = depth[u]+depth[v]-2*depth[lca]                      [5.4]
  |     +-- directions L/R/U: root-paths to their LCA, trim prefix     [5.4, LC 2096]
  |     +-- lca = min-depth on the 2n-1 tour window -> RMQ (Ch.11)     [5.4 preview]
  |
  +-- tree may be a CHAIN (n ~ 2e5) -> the stamp/tour recursion is O(h) deep:
      do it with an explicit (node,parent,stage) stack, or you RE.
      Full RMQ / binary-lifting LCA -> the LCA chapter [c11]; range structures ->
      [c13]/[c14]/[c15]; subtree-merge (sack) -> [c28]; offline blocks (Mo) -> [c21].
\end{lstlisting}

\chaptersection{Common Mistakes Checklist}

\begin{itemize}[nosep, leftmargin=1.4em]
  \item \textbf{Testing ancestry with the entry time alone.} $\tin(u)<\tin(v)$ says nothing --
        an earlier cousin passes it. Require $\tout(v)\le\tout(u)$ too (\secref{c24-times},
        Trap~1; \thmref{c24-anc}).
  \item \textbf{Forgetting the exit tick.} Without advancing the clock on the way out, every
        $\tout$ is stale and disjoint subtrees look nested (\secref{c24-times}, Trap~2).
  \item \textbf{Closed interval $[\tin,\tout]$ for a subtree.} It double-counts the next node;
        the range is half-open $[\tin,\tin{+}\size)$ (\secref{c24-flatten}, Trap~1;
        \thmref{c24-range}).
  \item \textbf{Indexing the flat array by node id.} Slot $k$ holds the node of preorder rank
        $k$, not node $k$ -- positions come from $\tin$ (\secref{c24-flatten}, Trap~3).
  \item \textbf{Re-DFS-ing per subtree query.} That is $O(\size[v])$ each, $O(n^2)$ overall --
        flatten once, answer by range (\secref{c24-flatten}, Trap~2).
  \item \textbf{Never undoing the path aggregate.} A value pushed on entry must pop on exit or
        sibling subtrees inherit a phantom prefix (\secref{c24-path}, Trap~1; \propref{c24-sweep}).
  \item \textbf{Reporting a $0$-based depth as the ``level.'} Problems index levels from $1$ --
        return \code{best+1} (\secref{c24-path}, Trap~2).
  \item \textbf{Building the tour with the entry-only clock.} The min-depth-LCA property needs
        the $2n{-}1$ walk that revisits on every return; a length-$n$ preorder has no re-visit
        window (\secref{c24-dist}, Trap~1; \defref{c24-tour}).
  \item \textbf{Sweeping the RMQ window un-ordered.} Order the endpoints with $\min$ and $\max$
        before scanning -- a raw window from $\mathrm{first}(u)$ to $\mathrm{first}(v)$ is empty
        (and wrong) when $u$ is met after $v$ (\secref{c24-dist}, Trap~2).
\end{itemize}

\chaptersection{WA Clinic}

Verdicts are symptoms; the causes are the short list this chapter drilled. Diagnose
top-down -- every leaf names the training problem.

\begin{figure}[H]
\centering
\begin{tikzpicture}[diagram,
    v/.style={boxnode, fill=primaryblue!12, draw=primaryblue, text width=22mm, align=center},
    b/.style={boxnode, fill=difforange!12, draw=difforange, text width=12mm},
    l/.style={boxnode, fill=successgreen!10, draw=successgreen, text width=58mm,
              font=\scriptsize, align=left}]
  \node[v] (verdict) {\textbf{Verdict on the judge}};
  \node[b, below left=8mm and 6mm of verdict] (wa) {\textbf{WA}};
  \node[b, below right=1mm of verdict] (tle) {\textbf{TLE}};
  \node[b, below left=8mm and 1mm of verdict, xshift=30mm] (re) {\textbf{RE}};
  \node[b, below right=8mm and 6mm of verdict] (mle) {\textbf{MLE}};
  \node[l, below right=1mm and 2mm of wa.south] (wa1)
    {ancestry from $\tin$ alone flags cousins (\probref{c24-throne}); subtree range closed by
     one grabs a neighbour (\probref{c24-fuel}); level reported $0$-indexed (\probref{c24-levelsum});
     RMQ window endpoints not min/max'd (\probref{c24-leafpairs})};
  \node[l, below right=1mm and 2mm of tle.south] (tle1)
    {DFS the subtree per query instead of reading one range (\probref{c24-fuel}); $O(n)$ naive
     LCA scan under a query load -- take a real RMQ/lifting structure (\chref{c11})};
  \node[l, below right=1mm and 2mm of re.south] (re1)
    {recursive stamp/tour on a $2\cdot10^{5}$ chain -- go explicit (node,parent,stage) stack
     (\secref{c24-times}, Trap~3; \secref{c24-flatten}, Trap~3)};
  \node[l, below right=1mm and 2mm of mle.south] (mle1)
    {materialising every root-path instead of a running sum (\secref{c24-path}); a $2n{-}1$ tour
     for a purely static subtree-sum problem that only needed $\tin/\size$ (\secref{c24-flatten})};
  \draw[flow] (verdict) -- (wa); \draw[flow] (verdict) -- (tle);
  \draw[flow] (verdict) -- (re); \draw[flow] (verdict) -- (mle);
  \draw[flow] (wa) -- (wa1); \draw[flow] (tle) -- (tle1);
  \draw[flow] (re) -- (re1); \draw[flow] (mle) -- (mle1);
\end{tikzpicture}
\caption*{Chapter~5 verdict tree -- only the failures this chapter's problems can produce.}
\end{figure}

\medskip
\noindent Backup table (same content, terse):
\begin{center}
\small
\renewcommand{\arraystretch}{1.2}
\begin{tabular}{@{}llp{0.42\linewidth}l@{}}
\toprule
\textbf{Verdict} & \textbf{Usual cause} & \textbf{One-line fix} & \textbf{Trained in} \\
\midrule
WA  & $\tin$-only ancestry test     & require $\tout(v)\le\tout(u)$ too          & \probref{c24-throne} \\
WA  & closed subtree interval        & half-open $[\tin,\tin{+}\size)$           & \probref{c24-fuel} \\
WA  & $0$-based ``level'' reported   & return \code{best+1}                         & \probref{c24-levelsum} \\
TLE & DFS per subtree query          & flatten once, range-read                  & \probref{c24-fuel} \\
RE  & recursion on a chain           & explicit (node,parent,stage) stack        & \secref{c24-times} \\
MLE & store every path / tour        & a running scalar or just $\tin/\size$     & \secref{c24-path} \\
\bottomrule
\end{tabular}
\end{center}

\chaptersection{Quick Reference Card}

\begin{quickref}[Chapter 5 at a glance]
\textbf{One question.} A subtree is a branching shape; \emph{when do I enter and leave each
node} turns it into a line. One DFS writes two integers per node and every tree question
becomes an array question (\defref{c24-times}, \thmref{c24-range}).

\medskip
\textbf{Two clocks, both yours.} A \emph{both-tick} clock ($\tin,\tout\in[0,2n)$) gives the
$O(1)$ ancestor test and the bracket view (\thmref{c24-anc}, \propref{c24-bracket}); an
\emph{entry-only} preorder rank ($\tin\in[0,n)$, $\tout=\tin+\size$) indexes the flat array
directly (\thmref{c24-range}). The two agree on containment -- pick by what you must index.

\medskip
\textbf{Reads of the same tour.} Subtree $\to$ range $[\tin,\tin{+}\size)$ (5.2). Root-path
$\to$ push/undo a running value down the DFS stack (5.3). Level $\to$ bucket by $\depth$ (5.3).
Distance $\to$ $\depth u+\depth v-2\depth{\lca}$; $\lca$ $\to$ min-depth on the $2n{-}1$ tour
between first-occurrences -- an RMQ you build in \chref{c11} (\thmref{c24-dist},
\propref{c24-rmqpreview}).
\end{quickref}

\begin{lstlisting}[caption={Chapter 5 code kit: stamp, flatten, path, tour}, label={lst:c24-kit}]
#include <bits/stdc++.h>
using namespace std;

vector<vector<int>> adj;
vector<int> tin_, tout_, sz_, dep_, flat_, pref_, sx_;
int timer_;

// (1) BOTH-tick clock: tin on entry, tout on exit -- O(1) ancestry. [5.1]
void stamp(int u, int p) {
    tin_[u] = timer_++;
    for (int v : adj[u]) if (v != p) stamp(v, u);
    tout_[u] = timer_++;
}
bool isAnc(int u, int v) {                       // u a proper ancestor of v
    return tin_[u] <= tin_[v] && tout_[v] <= tout_[u];
}

// (2) ENTRY-ONLY flatten: subtree(v) = flat range [tin, tin+size). [5.2]
void flatten(int u, int p, const vector<int>& w) {
    tin_[u] = timer_++; sz_[u] = 1; flat_[tin_[u]] = w[u];
    for (int v : adj[u]) if (v != p) { flatten(v, u, w); sz_[u] += sz_[v]; }
    tout_[u] = tin_[u] + sz_[u];
}
long long subtreeSum(int v) {                    // static, via prefix array
    return pref_[tout_[v]] - pref_[tin_[v]];
}

// (3) ROOT-PATH sweep: push on entry, UNDO on exit.                 [5.3]
long long pathSum = 0;
void path(int u, int p, const vector<int>& w) {
    pathSum += w[u];                             // value here = root->u sum
    for (int v : adj[u]) if (v != p) path(v, u, w);
    pathSum -= w[u];                             // the pop that trips everyone
}

// (4) 2n-1 TOUR for distance/LCA (RMQ proved in Ch.11).             [5.4]
vector<int> tour_, tourDep_, first_;
void tour(int u, int p, int d) {
    first_[u] = (int)tour_.size(); dep_[u] = d;
    tour_.push_back(u); tourDep_.push_back(d);
    for (int v : adj[u]) if (v != p) { tour(v, u, d + 1);
        tour_.push_back(u); tourDep_.push_back(d); }   // return step
}
\end{lstlisting}

\chaptersection{Write-From-Memory Drills}

Five signature skeletons, each lifted from a listing already in the chapter. Cover the key
line, reproduce it from memory, then check the answer one line below.

\begin{lstlisting}[caption={Drill 1 -- stamp entry/exit and test ancestry (from \coderef{c24-stamp})}]
void stamp(int u, int p) {
    tin_[u] = timer_++;
    for (int v : adj[u]) if (/* FILL: ______ */) stamp(v, u);
    tout_[u] = /* FILL: ______ */;
}
// isAnc(u,v): return tin_[u] <= tin_[v] && /* FILL: ______ */;
\end{lstlisting}
\noindent{\scriptsize Key: parent guard \code{v != p}; the exit tick \code{tout[u]=timer++};
and the containment needs \emph{both} endpoints \code{tout[v] <= tout[u]} (\thmref{c24-anc}).}

\begin{lstlisting}[caption={Drill 2 -- flatten so a subtree is a range (from \coderef{c24-flatten-build})}]
void flatten(int u, int p, const vector<int>& w) {
    tin_[u] = timer_++; sz_[u] = 1; flat_[/* FILL: __ */] = w[u];
    for (int v : adj[u]) if (v != p) { flatten(v, u, w); sz_[u] += sz_[v]; }
    tout_[u] = /* FILL: tin_[u] + sz_[u] */;
}
\end{lstlisting}
\noindent{\scriptsize Key: place the value at slot $\tin$, not at node id; the exclusive end
is $\tin+\size$ (\thmref{c24-range}). This is Trap~3 of \secref{c24-flatten}.}

\begin{lstlisting}[caption={Drill 3 -- subtree sum by prefix read (from \coderef{c24-flatten-build})}]
long long subtreeSum(int v) {
    int l = /* FILL: ______ */, r = /* FILL: ______ */;
    return pref_[r] - pref_[l];
}
\end{lstlisting}
\noindent{\scriptsize Key: $l=\tin[v]$, $r=\tin[v]+\size[v]$ (half-open); a static range sum is
one subtraction of prefix values (\propref{c24-prefix}).}

\begin{lstlisting}[caption={Drill 4 -- root-path aggregate with undo (from \coderef{c24-path-stack})}]
void path(int u, int p, const vector<int>& w) {
    pathSum += w[u];
    for (int v : adj[u]) if (v != p) path(v, u, w);
    /* FILL: ______ */;                  // <-- do not forget this line
}
\end{lstlisting}
\noindent{\scriptsize Key: \code{pathSum -= w[u]} restores the sibling branch; the value seen
\emph{before} recursing is the root-to-$u$ sum (\propref{c24-sweep}, Trap~1).}

\begin{lstlisting}[caption={Drill 5 -- build the 2n-1 tour (from \coderef{c24-tour-build})}]
void tour(int u, int p, int d) {
    first_[u] = (int)tour_.size();
    tour_.push_back(u); tourDep_.push_back(d);
    for (int v : adj[u]) if (v != p) { tour(v, u, d + 1);
        /* FILL: push u again */;         // the return step makes length 2n-1
    }
}
\end{lstlisting}
\noindent{\scriptsize Key: after each child, record $u$ again -- the LCA window only exists in
the revisiting tour (\defref{c24-tour}, Trap~1).}

\chaptersection{Complete Problem Checklist}

Every problem this chapter offers -- 6 fully solved in-section, 4 practice, one Summit, five
mixed. Tick \emph{only} after an unaided solve; a hinted attempt counts as half. Cite by name
and number -- URLs appear once, in the solved and practice entries.

\begin{center}
\footnotesize
\renewcommand{\arraystretch}{1.25}
\begin{tabular}{@{}c p{0.17\linewidth} l p{0.20\linewidth} p{0.24\linewidth} c@{}}
\toprule
\textbf{\#} & \textbf{Problem} & \textbf{Source} & \textbf{Key technique} & \textbf{What to look for} & \textbf{Done?} \\
\midrule
\multicolumn{6}{l}{\textbf{Section 5.1 -- times and bracket sequences}} \\
W5.1.1 & Throne Inheritance & LC 1600 & \lv3 preorder $=$ succession & list survivors in $\tin$ order, filter dead & $\square$ \\
P5.1.2 & Binary Tree Coloring Game & LC 1145 & \lv3 subtree size span & compare $\size[v]$ to $n-\size[v]$ & $\square$ \\
\midrule
\multicolumn{6}{l}{\textbf{Section 5.2 -- subtree as interval}} \\
W5.2.1 & Minimum Fuel Cost to Capital & LC 2477 & \lv3 subtree size per edge & ceil(span/capacity) over every edge, \code{long long} & $\square$ \\
P5.2.2 & Reachable Nodes w/ Restrictions & LC 2368 & \lv3 skip a subtree & $n$ minus the disjoint restricted spans & $\square$ \\
\midrule
\multicolumn{6}{l}{\textbf{Section 5.3 -- path and depth sequence}} \\
W5.3.1 & Maximum Level Sum & LC 1161 & \lv2 level buckets & per-depth sum, argmax, 1-index, seed $-\infty$ & $\square$ \\
W5.3.2 & Kth Largest Level Sum & LC 2583 & \lv3 level $+$ order statistic & collect level sums, sort, guard $k$ & $\square$ \\
P5.3.3 & Linked List in Binary Tree & LC 1367 & \lv3 path-stack DFS & match list down a root-path, undo on fail & $\square$ \\
\midrule
\multicolumn{6}{l}{\textbf{Section 5.4 -- distance and the RMQ preview}} \\
W5.4.1 & Step-By-Step Directions & LC 2096 & \lv3 root-paths $+$ LCA & trim common prefix, emit U/L/R & $\square$ \\
W5.4.2 & Number of Good Leaf Pairs & LC 1530 & \lv4 cross-count at LCA & pair leaves from \emph{different} child lists & $\square$ \\
P5.4.3 & Diameter of N-Ary Tree & LC 1522 & \lv3 two tallest chains & best path through $u$ $=$ sum of top-2 heights & $\square$ \\
\midrule
\multicolumn{6}{l}{\textbf{Summit (chapter end)}} \\
S5 & Minimum Score After Removals & LC 2322 & \lv5 subtree XOR $+$ interval split & ancestor/disjoint case split, minimise $\max-\min$ & $\square$ \\
\bottomrule
\end{tabular}
\end{center}

\chaptersection{Recommended Study Order}

Five focused days. Days 1--2 fix the two clocks and the flattening; day 3 the path/level
reads; day 4 the tour and distance; day 5 the Summit. Solve in this order even where it
differs from section order -- each day's tool is already in hand the day before.

\begin{center}
\small
\renewcommand{\arraystretch}{1.25}
\begin{tabular}{@{}llp{0.55\linewidth}@{}}
\toprule
\textbf{Day} & \textbf{Focus} & \textbf{Tasks} \\
\midrule
1 & Read + do 5.1 & stamp $G$ by hand (\exref{c24-times-run}); LC 1600; write \code{isAnc} cold \\
2 & Read + do 5.2 & flatten $G$ (\exref{c24-flat-run}); LC 2477; LC 2368; LC 1145 \\
3 & Read + do 5.3 & path push/undo; LC 1161, LC 2583; LC 1367 \\
4 & Read + do 5.4 & build the $2n{-}1$ tour (\exref{c24-dist-run}); LC 2096, LC 1530; LC 1522 \\
5 & Summit + revise & LC 2322 subtree-XOR; walk the decision flowchart on five questions \\
\bottomrule
\end{tabular}
\end{center}

\chaptersection{Difficulty Progression}

Tiers are competency gates -- each row names the \emph{skill} the tier unlocks, not just a
rating bucket.

\begin{center}
\small
\renewcommand{\arraystretch}{1.3}
\begin{tabular}{@{}lp{0.36\linewidth}p{0.34\linewidth}@{}}
\toprule
\textbf{Tier} & \textbf{Problems} & \textbf{What passing this tier proves} \\
\midrule
$\lv2$ \ $\bigstar\bigstar$
  & LC 1161
  & you can bucket nodes by depth and read a per-level answer off one array \\
$\lv3$ \ $\bigstar\bigstar\bigstar$
  & LC 1600, 1145, 2477, 2368, 2583, 1367, 2096, 1522
  & one DFS's $\tin/\tout/\size$ answers ancestry, subtree size, and root-path questions \\
$\lv4$ \ $\bigstar\bigstar\bigstar\bigstar$
  & LC 1530
  & combine two child-side aggregates at a common ancestor (a distance/pair count) \\
$\lv5$ \ $\bigstar\bigstar\bigstar\bigstar\bigstar$
  & LC 2322 (Summit)
  & reduce edge-deletion to interval splits and case-split on ancestor containment \\
\bottomrule
\end{tabular}
\end{center}

\chaptersection{Reduction Gallery}

The tour is a disguise; each line collapses a family onto one reading of it. No new theory --
just where $\tin/\tout$ is already hiding.

\begin{itemize}[nosep, leftmargin=1.4em]
  \item ``is $u$ an ancestor of $v$? / the $k$-th successor'' $\longrightarrow$ the $O(1)$
        containment test on $[\tin,\tout)$ (\thmref{c24-anc}).
  \item ``a line of succession in DFS order'' (LC 1600) $\longrightarrow$ list nodes by $\tin$,
        filter (\probref{c24-throne}).
  \item ``sum/count/max over the whole subtree of $v$'' (LC 2368, 2477) $\longrightarrow$ the
        flat range $[\tin,\tin{+}\size)$ (\thmref{c24-range}).
  \item ``cars forced over each edge'' / per-edge subtree counts (LC 2477) $\longrightarrow$ the
        span $\tout-\tin=\size$ of the child subtree (\probref{c24-fuel}).
  \item ``sum on the path from the root to $v$'' $\longrightarrow$ push/undo a running value on
        the DFS stack (\propref{c24-sweep}).
  \item ``an answer per level'' (LC 1161, 2583) $\longrightarrow$ the depth sequence: bucket by
        $\depth$, then scan/sort an ordinary array (\probref{c24-levelsum}, \probref{c24-kthlevel}).
  \item ``directions / distance between two nodes'' (LC 2096) $\longrightarrow$ $\dist$ via the
        LCA, or stitch two root-paths to their common prefix (\thmref{c24-dist},
        \probref{c24-pathdir}).
  \item ``count leaf pairs within distance $d$'' (LC 1530) $\longrightarrow$ merge child depth
        lists at their common ancestor (\probref{c24-leafpairs}).
  \item ``longest path in the tree'' (LC 1522) $\longrightarrow$ top-2 downward heights at each
        node -- a distance whose LCA is the node itself (\secref{c24-dist}).
  \item ``cut two edges, score three pieces'' (LC 2322) $\longrightarrow$ each cut is a subtree
        interval; the pieces are subtree XORs combined by ancestor containment (\probref{c24-summit}).
\end{itemize}

\chaptersection{Mixed Problem Set}

These refuse to name their pattern; decide \emph{which reading of the tour} you need before
opening the code. A hint tells you the shape of the answer, never the tool.

\begin{mixedset}
  \practice{Range Sum of a BST}{LeetCode 938}
    {https://leetcode.com/problems/range-sum-of-bst/}{LC Easy L2}
    {the answer is a sum over the nodes whose \emph{values} lie in $[\text{low},\text{high}]$
    -- a range query, but on the value axis rather than $\tin$. Prune a subtree the moment it
    falls outside the interval: the BST order makes each skipped side one contiguous block,
    the same ``a subtree is a range'' economy as \secref{c24-flatten}.}
  \practice{Find Nearest Right Node in Binary Tree}{LeetCode 1779}
    {https://leetcode.com/problems/find-nearest-right-node-in-binary-tree/}{LC Easy L2}
    {``nearest node at the same depth, to the right'' is a question about the \emph{depth
    sequence} -- bucket nodes by $\depth$ (\secref{c24-path}) in left-to-right order, then read
    the successor of the target inside its own level's array.}
  \practice{Maximum Sum BST in Binary Tree}{LeetCode 1373}
    {https://leetcode.com/problems/maximum-sum-bst-in-binary-tree/}{LC Medium L4}
    {a subtree's validity and sum are decided \emph{bottom-up} from each child's summary
    (min, max, sum, is-BST) -- the subtree-as-a-unit thinking of \secref{c24-flatten}, but the
    aggregate is a fixed-size tuple instead of a range sum. Combine two children at their
    parent, exactly as \probref{c24-leafpairs} pairs across an ancestor.}
  \practice{Trim a Binary Search Tree}{LeetCode 669}
    {https://leetcode.com/problems/trim-a-binary-search-tree/}{LC Easy L2}
    {keeping only values in $[\text{low},\text{high}]$ means \emph{dropping whole subtrees} that
    fall outside -- recognise that a BST node with value below $\text{low}$ carries an entire
    discarded left interval. Recurse on the surviving range and reattach, the pointer cousin of
    the ``paint/keep a contiguous subtree'' idea in \secref{c24-flatten}.}
  \practice{Closest Nodes Queries in a Binary Search Tree}{LeetCode 2374}
    {https://leetcode.com/problems/closest-nodes-queries-in-a-binary-search-tree/}{LC Medium L3}
    {many point queries over one ordered tree: flatten the BST to its \emph{inorder} array
    (an ordered tour of the nodes), then each query is a $\log$ binary search / two-neighbour
    probe on that sorted line -- the flatten-to-a-line move of \secref{c24-flatten} applied to
    a sorted sequence.}
  \practice{Subtree Queries}{CSES 1137}
    {https://cses.fi/problemset/task/1137}{CSES · L4}
    {Node values, two query kinds: add to a whole subtree, or read one node. Lay the
    tree on a line once with \secref{c24-flatten} and both become array work --
    \thmref{c24-range}'s "a subtree is one contiguous block" is the entire
    translation; the rest is a difference count or a Fenwick pair.}
  \practice{Distinct Values Queries}{CSES 1734}
    {https://cses.fi/problemset/task/1734}{CSES · L4}
    {How many distinct values lie in a range $[a,b]$. Sort the queries by right end and sweep,
    keeping only the \emph{latest} occurrence of each value alive in a Fenwick -- so a range
    count is one prefix difference. (Turn a subtree-distinct question into this exact array
    problem by first laying the tree on a line with \secref{c24-flatten}.)}
  \practice{Fixed-Length Paths I}{CSES 2080}
    {https://cses.fi/problemset/task/2080}{CSES · L5}
    {Count the node pairs at distance exactly $k$. Root once and keep the depth
    buckets of the part of the tree already explored as the DFS descends: a new node
    partners with the marked ones $k$ levels up -- \probref{c24-leafpairs}'s
    "combine lists at the common ancestor" with buckets instead of merges.}
  \practice{Maximum Points After Collecting Coins From All Nodes}{LeetCode 2920}
    {https://leetcode.com/problems/maximum-points-after-collecting-coins-from-all-nodes/}{LC Hard L5}
    {A walk may reverse direction once on the way back to the root -- state that as a
    fixed-size table per subtree and combine bottom-up: the rooted-line thinking of
    \secref{c24-flatten}, with a DP table where the chapter kept sums.}
  \practice{Find Minimum Diameter After Merging Two Trees}{LeetCode 3203}
    {https://leetcode.com/problems/find-minimum-diameter-after-merging-two-trees/}{LC Hard L4}
    {One bridge between the trees, then the merged diameter must be as small as
    possible: only the \emph{centres} are worth attaching, because the radius is the
    smallest maximum distance any root-path machinery can offer -- two BFS sweeps per
    tree, the distance thinking of \thmref{c24-dist} used as a measuring stick.}
\end{mixedset}

\chaptersection{The Summit -- LC 2322 Minimum Score After Removals on a Tree}

\begin{notebox}[Why this is the chapter's top rope]
Every earlier section read \emph{one} tree through \emph{one} interval: an ancestry test, a
subtree range, a path, a distance. The Summit makes you delete \emph{two} edges at once and
score the three resulting pieces -- and the only clean way to know which nodes end up in which
piece is to have laid the tree on a line. Each cut edge is a child node's interval; a second
cut either falls \emph{inside} the first interval (nested) or in a disjoint one, and the three
component XORs are just subtree XORs combined differently by that one ancestor test. It is the
exact composition the chapter kept promising -- the $O(1)$ containment of
\thmref{c24-anc} and the subtree block of \thmref{c24-range} -- doing the whole bookkeeping for
an $O(n^2)$ scan over edge pairs on $n\le 1000$.
\end{notebox}

\solvedproblem{c24-summit}{Minimum Score After Removals on a Tree}{LeetCode 2322}
  {https://leetcode.com/problems/minimum-score-after-removals-on-a-tree/}{LC Hard L5}
  {root it, subtree XOR $\mathit{sx}[u]$, then $\tin/\tout$ ancestor containment to case-split
   every edge pair into nested or disjoint (\thmref{c24-anc}, \thmref{c24-range})}

\begin{problemstatement}[Minimum Score After Removals on a Tree]
An undirected connected \textbf{tree} with $n$ nodes labelled $0\dots n{-}1$ ($n$ even,
$4\le n\le 1000$) is given by a node-value array \code{nums} and $n-1$ edges
\code{edges[i]=[a,b]}. Remove \textbf{two distinct edges}, splitting the tree into exactly
\textbf{three} connected components. Let the XOR of all node values in a component be its XOR;
the \term{score} of a removal is $\max(\text{three XORs})-\min(\text{three XORs})$. Return the
\textbf{minimum} score over all valid pairs of removed edges.
\end{problemstatement}

\begin{attemptgate}[Attempt it yourself first -- \lv{5}, give yourself 45 minutes]
A fair first attempt: enumerate all $\binom{n-1}{2}$ edge pairs, run a BFS/DFS on the remaining
graph to read off the three components, XOR each. Correct, but re-walking the forest per pair is
$\bigO{n^3}$ -- the challenge is computing every component XOR \emph{from timestamps}, not from
a fresh traversal.
\peek{root the tree once. What does deleting the edge above node $u$ ``remove'' -- in terms of
one array you can precompute during the DFS?}
\nudge{the subtree of $u$: its XOR is $\mathit{sx}[u]$, and the rest of the tree is
$s\oplus\mathit{sx}[u]$ where $s$ is the total XOR ($a\oplus a=0$). Now, for \emph{two} cut nodes
$a,b$, is one inside the other's subtree? That is exactly the containment test of
\thmref{c24-anc}.}
\fullhint{Precompute $\tin/\tout$ and subtree XOR $\mathit{sx}[\cdot]$ in one DFS. Each non-root
node is ``the edge to its parent.'' For every pair $(a,b)$: if $a$ is an ancestor of $b$ the
pieces are $\mathit{sx}[b]$, $\mathit{sx}[a]\oplus\mathit{sx}[b]$, $s\oplus\mathit{sx}[a]$;
if disjoint they are $\mathit{sx}[a]$, $\mathit{sx}[b]$, $s\oplus\mathit{sx}[a]\oplus\mathit{sx}[b]$
(the symmetric case swaps $a,b$). Score $=\max-\min$; take the minimum. $O(n^2)$. Full build and
proof below.}
\end{attemptgate}

\solstep{Understanding the Problem}
Deleting two edges creates three pieces, but \emph{which} nodes are in which piece depends on
where the two cuts sit in the rooted tree. Once rooted at $0$, every edge is ``above some child
node $u$,'' and cutting it peels off exactly $\subtree(u)$. Two cuts are therefore either
\emph{nested} (one subtree inside the other) or \emph{disjoint} -- and the three piece XORs are
combinations of the two subtree XORs and the total. The whole problem becomes: precompute the
subtree XORs, decide nesting in $O(1)$, and try all $\binom{n-1}{2}$ pairs.

\solstep{Key Observations}
\begin{observation}
XOR cancels itself: $a\oplus a = 0$. With total $s=\bigoplus_v \text{nums}[v]$ and subtree
XOR $\mathit{sx}[u]$, the complement of $\subtree(u)$ is $s\oplus \mathit{sx}[u]$ -- no second
traversal needed.
\end{observation}
\begin{observation}
$\mathit{sx}[u]=\text{nums}[u]\oplus\bigoplus_{c\to u}\mathit{sx}[c]$ is a bottom-up return of
one DFS, computed alongside $\tin/\tout$ in the same pass (\thmref{c24-range}'s clock).
\end{observation}
\begin{observation}
For two cut nodes $a,b$: $a$ is an ancestor of $b$ iff $\mathit{tin}[a]\le\mathit{tin}[b]$ and
$\mathit{tout}[b]\le\mathit{tout}[a]$ (\thmref{c24-anc}). This single test picks the case split.
\end{observation}

\solstep{Think It Through With Me}
\begin{thinkbox}
\deadend{For each edge pair, delete them and DFS the forest for the three components.}
{$O(n)$ per pair times $O(n^2)$ pairs is $\bigO{n^3}$ -- about $10^{9}$ at $n=1000$, a TLE. The
three components are not arbitrary: they are always $\subtree(\cdot)$ pieces, so their XORs are
already in $\mathit{sx}[\cdot]$.}
\deadend{Treat every pair with one formula.}{A nested pair and a disjoint pair give
\emph{different} three pieces. Writing them as one expression double-counts or drops a subtree --
the two cases must be separated, and \thmref{c24-anc} is the switch that separates them.}
\aha{Store $\mathit{sx},\tin,\tout$ and the total $s$. For each ordered pair of cut nodes, test
containment, apply the matching three-XOR combination, score $\max-\min$, and keep the minimum.
Two linear DFS passes plus an $O(n^2)$ scan.}
\end{thinkbox}

\solstep{The Approach Ladder}
\begin{center}\renewcommand{\arraystretch}{1.3}
\begin{tabular}{@{}c p{0.235\linewidth} p{0.30\linewidth} l c@{}}
\toprule
\textbf{\#} & \textbf{Approach} & \textbf{Idea} & \textbf{Time} & \textbf{Verdict} \\
\midrule
1 & \lv{3} re-DFS per edge pair   & delete two edges, BFS three pieces  & $O(n^3)$ & \tle \\
2 & \lv{4} subtree sums, one case & nested only, forget disjoint        & $O(n^2)$ & \textbf{WA} \\
3 & \lv{5} subtree XOR + containment split & $\mathit{sx}[\cdot]$ + \thmref{c24-anc} case split & $O(n^2)$ & \ac \\
\bottomrule
\end{tabular}\end{center}
Rung 1 is honest but cubic. Rung 2 precomputes the pieces but mis-scores the disjoint case (the
classic WA). Rung 3 keeps the precomputation and adds the single ancestor test that tells the
two geometries apart.

\solstep{Solution}
Root at $0$; one DFS fills $\tin/\tout$ (both- or entry-only clock -- the containment test works
for either, here the entry-only clock with $\tout=\tin+\size$), $\size$, and the subtree XOR
$\mathit{sx}$. Enumerate unordered node pairs $(a,b)$ with $a<b$, both non-root: each is the edge
to its parent. Apply the nested/disjoint case split to get three XORs, score them
$\max-\min$, and take the global minimum (\coderef{c24-summit}).

\solstep{Proof of Correctness}
\begin{proofbox}[Two lemmas, then a case split]
\emph{Lemma 1 (pieces from subtree XORs).} Deleting the edge above node $u$ splits the tree into
$\subtree(u)$, with XOR $\mathit{sx}[u]$, and its complement, with XOR $s\oplus\mathit{sx}[u]$
because $s=\mathit{sx}[u]\oplus(s\oplus\mathit{sx}[u])$ and $\oplus$ is its own inverse.
\emph{Lemma 2 (the second cut refines one piece).} Given two cut nodes $a,b$: if $a$ is an
ancestor of $b$, then $\subtree(b)\subseteq\subtree(a)$, so the three components are
$\subtree(b)$, $\subtree(a)\setminus\subtree(b)$, and the outside -- XORing to
$\mathit{sx}[b]$, $\mathit{sx}[a]\oplus\mathit{sx}[b]$, and $s\oplus\mathit{sx}[a]$; the symmetric
statement holds when $b$ is an ancestor of $a$; if neither is an ancestor of the other,
$\subtree(a)$ and $\subtree(b)$ are disjoint and the components are $\subtree(a)$,
$\subtree(b)$, and the rest, XORing to $\mathit{sx}[a]$, $\mathit{sx}[b]$,
$s\oplus\mathit{sx}[a]\oplus\mathit{sx}[b]$. \emph{Exhaustion.} Every valid removal is a pair of
non-root nodes, exactly one of these three geometries holds (decided by \thmref{c24-anc}, whose
correctness is the interval-containment characterisation), and the score is
$\max-\min$ of the three XORs. Scanning all $\binom{n-1}{2}$ pairs and returning the least score
therefore returns the minimum possible score.
\end{proofbox}

\solstep{Implementation}
\begin{lstlisting}[caption={Minimum Score After Removals on a Tree (LeetCode 2322) -- subtree XOR + tin/tout case split}, label={lst:c24-summit}]
#include <bits/stdc++.h>
using namespace std;

class Solution {
    vector<vector<int>> adj;
    vector<int> val_, tin_, tout_, sz_, sx_;
    int timer_, total_;

    void dfs(int u, int p) {                 // entry-only clock; tout = tin + size
        tin_[u] = timer_++;
        sz_[u]  = 1;
        sx_[u]  = val_[u];
        for (int v : adj[u]) if (v != p) {
            dfs(v, u);
            sz_[u] += sz_[v];
            sx_[u] ^= sx_[v];                 // fold in each child subtree XOR
        }
        tout_[u] = tin_[u] + sz_[u];
    }
    bool isAnc(int u, int v) {               // u a proper ancestor of v (thm:c24-anc)
        return u != v && tin_[u] <= tin_[v] && tout_[v] <= tout_[u];
    }
    int scoreOf(int a, int b) {              // cut the parent-edges of nodes a and b
        int x, y, z;
        if (isAnc(a, b))       { x = sx_[b];        y = sx_[a] ^ sx_[b]; z = total_ ^ sx_[a]; }
        else if (isAnc(b, a))  { x = sx_[a];        y = sx_[b] ^ sx_[a]; z = total_ ^ sx_[b]; }
        else                   { x = sx_[a];        y = sx_[b];           z = total_ ^ sx_[a] ^ sx_[b]; }
        return max({x, y, z}) - min({x, y, z});
    }
public:
    int minimumScore(vector<int>& nums, vector<vector<int>>& edges) {
        int n = nums.size();
        adj.assign(n, {}); val_ = nums;
        tin_.assign(n, 0); tout_.assign(n, 0); sz_.assign(n, 0); sx_.assign(n, 0);
        for (auto& e : edges) { adj[e[0]].push_back(e[1]); adj[e[1]].push_back(e[0]); }
        total_ = 0; for (int v : nums) total_ ^= v;
        timer_ = 0; dfs(0, -1);
        int ans = INT_MAX;
        for (int a = 1; a < n; ++a)          // node a != root  ==  the edge (parent[a], a)
            for (int b = a + 1; b < n; ++b)
                ans = min(ans, scoreOf(a, b));
        return ans;
    }
};
\end{lstlisting}

\solstep{Trace}
\begin{tracebox}[Worked example -- a $6$-node tree]
\emph{Input:} \code{nums = [2,1,3,4,5,6]}, \code{edges = [[0,1],[0,2],[1,3],[1,4],[4,5]]}.
Root $0$; children $0\!\to\!\{1,2\}$, $1\!\to\!\{3,4\}$, $4\!\to\!\{5\}$; preorder
$0,1,3,4,5,2$. Then $\mathit{tin}=[0,1,5,2,3,4]$, $\mathit{sz}=[6,4,1,1,2,1]$,
$\mathit{tout}=[6,5,6,3,5,5]$. Subtree XORs bottom-up:
$\mathit{sx}=[7,6,3,4,3,6]$ and the total $s=2\oplus1\oplus3\oplus4\oplus5\oplus6=7$
($=\mathit{sx}[0]$, a free check). Cut nodes are $\{1,2,3,4,5\}$. Representative pairs:
\begin{itemize}[nosep,leftmargin=1.3em]
  \item $(1,2)$ disjoint: pieces $\mathit{sx}[1]{=}6$, $\mathit{sx}[2]{=}3$,
        $s\oplus6\oplus3{=}2$; score $6-2=4$.
  \item $(1,3)$ nested ($1$ anc of $3$): $\mathit{sx}[3]{=}4$, $\mathit{sx}[1]\oplus\mathit{sx}[3]{=}2$,
        $s\oplus\mathit{sx}[1]{=}1$; score $4-1=3$.
  \item $(4,5)$ nested ($4$ anc of $5$): $\mathit{sx}[5]{=}6$, $\mathit{sx}[4]\oplus\mathit{sx}[5]{=}5$,
        $s\oplus\mathit{sx}[4]{=}4$; score $6-4=\mathbf{2}$.
\end{itemize}
Over all $10$ pairs the least score is $2$, attained by deleting edges $(1,4)$ and $(4,5)$
-- the components $\{5\}$ (XOR $6$), $\{4\}$ (XOR $5$), $\{0,1,2,3\}$ (XOR $2\oplus1\oplus3\oplus4=4$).
Answer $=2$. \checkmark
\end{tracebox}

\solstep{Complexity}
\begin{complexitybox}
Time $\bigO{n^2}$ -- one $\bigO{n}$ DFS to build $\tin/\tout/\size/\mathit{sx}$, then
$\binom{n-1}{2}$ pairs each scored in $O(1)$ by the containment test. Space $\bigO{n}$ for the
five arrays and adjacency list, plus $\bigO{h}$ recursion. At $n=1000$ that is about
$5\cdot10^{5}$ pair scores -- comfortably inside limits. On a chain-shaped tree the DFS
recursion is $\bigO{n}$ deep; $n\le1000$ keeps it safe, but see the hardening box for the
iterative rewrite.
\end{complexitybox}

\solstep{Verification and Edge Cases}
\begin{center}\renewcommand{\arraystretch}{1.2}
\begin{tabular}{@{}p{0.30\linewidth} p{0.62\linewidth}@{}}
\toprule
\textbf{Case} & \textbf{What the code does} \\
\midrule
Smallest $n=4$ & $3$ non-root nodes, $\binom{3}{2}=3$ pairs; every pair is nested or disjoint \\
Best pair is disjoint & handled by the \code{else} branch (two separate subtrees) \\
Best pair is nested & handled by \code{isAnc(a,b)} or its mirror \code{isAnc(b,a)} \\
Total XOR zero & still correct -- pieces are XOR combinations, no division involved \\
Two cuts on a single spine & the nested branch fires repeatedly; $\max/\min$ over three XORs \\
\bottomrule
\end{tabular}\end{center}

\solstep{Pitfalls and Variants}
\begin{warningbox}
Three classic bugs: (i) using the \emph{entry-only} $\tin$ but a $\tout$ that is a
\emph{separate} exit tick (both-tick clock) -- the containment test needs the two clocks to agree
(\thmref{c24-anc}); this listing uses $\tout=\tin+\size$ consistently. (ii) Forgetting the
\texttt{isAnc(b,a)} mirror, so a $b$-over-$a$ nesting falls into the disjoint branch and
mis-scores (that is exactly Rung~2's WA). (iii) Scoring $\max-\min$ on node \emph{counts} instead
of \emph{XORs}. Variant -- the identical skeleton answers the ``remove two edges, minimise the
spread of component \emph{sizes}'' version by swapping $\oplus$ for $+$; and this is the same
interval-containment bookkeeping that powers Mo-on-trees (\chref{c21}) and DSU-on-tree (\chref{c28}).
\end{warningbox}

\begin{insightbox}[Key takeaways -- three things to remember]
\begin{enumerate}[nosep, leftmargin=1.6em]
  \item An edge cut is a subtree \emph{interval}; two cuts are either \emph{nested} or
        \emph{disjoint}, and \thmref{c24-anc} tells you which in $O(1)$.
  \item The complement trick $s\oplus\mathit{sx}[u]$ (or $s-\mathit{sx}$ for sums) means you
        never re-traverse -- the pieces come out of one precomputed array (\thmref{c24-range}).
  \item The whole Hard reduces to the chapter's two simplest tools: an ancestor test and a
        subtree sum. Depth of insight, not depth of code.
\end{enumerate}
\end{insightbox}

\begin{hardeningbox}
\begin{itemize}[nosep, leftmargin=1.4em]
  \item \textbf{RE on a deep tree}: replace the recursive DFS with an explicit
        \code{(node,parent,stage)} stack -- build a preorder order iteratively, then walk it in
        \emph{reverse} to accumulate $\mathit{sx}$ and $\size$. Same $O(n)$, zero call frames.
  \item \textbf{Arbitrary component aggregate}: swap $\oplus$ for $+,\min,\max,\gcd$ -- whenever
        the piece value is a \emph{subtree reducible} aggregate, the nested/disjoint case split
        is unchanged; that catalogue is Tree DP (\chref{c10}) and the query structures
        (\chref{c13}--\chref{c15}).
  \item \textbf{Larger $n$}: $O(n^2)$ pair enumeration does not scale past $\sim10^4$; problems
        that ask for \emph{one} best cut or DSU-style merging point to DSU-on-tree
        (\chref{c28}) or rerooting (\chref{c20}) rather than this scan.
\end{itemize}
\end{hardeningbox}

\chaptersection{Tool Kit Inventory}

The machine assembled so far -- cumulative through Chapter~5. Chapters~1--4's rows are
carried forward; Chapter~5's rows are appended. Pointers only, never prose.

\begin{center}
\small
\renewcommand{\arraystretch}{1.25}
\begin{tabular}{@{}p{0.24\linewidth} l c p{0.27\linewidth} p{0.16\linewidth}@{}}
\toprule
\textbf{Tool} & \textbf{Since} & \textbf{Level} & \textbf{Trigger phrase} & \textbf{Where} \\
\midrule
\multicolumn{5}{l}{\itshape Chapters 1--4 -- carried forward} \\
Adjacency list + parent-guard DFS & 1.1 & \lv1 & edge list given, not pointers & \coderef{c01-build} \\
Top-down carry / bottom-up return & 1.2 & \lv1 & value from root / subtree aggregate & \coderef{c01-measure} \\
Structural skeleton (null base + 2 calls) & 2.1 & \lv1 & any node-plus-two-children question & \coderef{c02-skeleton} \\
Return-new-root transform & 2.4 & \lv3 & the pointer graph itself must change & \coderef{c02-transform-skeleton} \\
Three-slot DFS (pre/in/post) & 3.1 & \lv1 & a node sequence in a named order & \coderef{c03-dfs} \\
Explicit-stack traversal / Morris & 3.2 & \lv2--\lv4 & recursion depth under attack & \coderef{c03-inorder-stack} \\
Level-order BFS with frontier snapshot & 3.3 & \lv2 & output or reduce \emph{per depth} & \coderef{c03-bfs} \\
Four storage models + children fetch & 4.1 & \lv1 & commit edges to memory; reach kids & \coderef{c04-reps} \\
Root: read parent, depth, size, height & 4.2 & \lv3 & undirected edges, need above/below & \coderef{c04-root} \\
Reroot shift $f(c)=f(p)+(n-2\size)$ & 4 (Summit) & \lv5 & the same aggregate from \emph{every} root & \coderef{c04-summit} \\
\midrule
\multicolumn{5}{l}{\itshape Chapter 5 -- new this chapter} \\
Stamp $\tin/\tout$; $O(1)$ ancestor test & 5.1 & \lv2 & ``is $u$ above $v$?'' / DFS-order listing & \coderef{c24-stamp} \\
Flatten: subtree $=$ range $[\tin,\tin{+}\size)$ & 5.2 & \lv3 & many whole-subtree queries & \coderef{c24-flatten-build} \\
Root-path push/undo sweep & 5.3 & \lv3 & aggregate along a root-to-node path & \coderef{c24-path-stack} \\
$2n{-}1$ tour $+$ distance via LCA & 5.4 & \lv4 & distance/directions between two nodes & \coderef{c24-tour-build} \\
Subtree XOR $+$ containment case split & Summit & \lv5 & ``cut two edges, score three pieces'' & \coderef{c24-summit} \\
\bottomrule
\end{tabular}
\end{center}

\chaptersection{Related Topics}

\begin{itemize}[nosep, leftmargin=1.4em]
  \item \chref{c03}, \chref{c04} -- backward: the DFS orders (\chref{c03}) and the
        parent-guarded rooting with $\parent/\depth/\size$ (\secref{c04-rooting}) are exactly
        what this chapter \emph{stamps}; nothing here re-derives them.
  \item \chref{c11} -- forward: the $2n{-}1$ tour's LCA is \emph{proved} and made
        fast (binary lifting) there; this chapter only previews the RMQ reduction
        (\propref{c24-rmqpreview}).
  \item \chref{c13}, \chref{c14}, \chref{c15} -- forward: the range \emph{answering/updating}
        over the flat array we hand off -- Fenwick for point-update, segment tree and lazy for
        range-add.
  \item \chref{c16} -- forward: Heavy--Light turns \emph{any} path (not just root-paths) into
        ranges; this chapter's 5.3 is its root-only special case.
  \item \chref{c20}, \chref{c28}, \chref{c21}, \chref{c22} -- the neighbouring subtree
        techniques: rerooting, DSU-on-tree/sack, Mo-on-trees, and hashing all \emph{start} from
        the $\tin/\tout$ built here.
\end{itemize}

\begin{lastword}[Last Word -- Chapter 5]
\emph{``A tree is recursion wearing a branching costume; an Euler tour is that costume taken
off, revealing one long array underneath.''}

\medskip
Chapters~3 and~4 taught you to \emph{walk} a tree and to \emph{store} it. This chapter asked the
question that turns a tree into something a segment tree can eat: \emph{when do I enter and leave
each node}. Two integers per node -- a birth time and a death time -- and the whole branching
shape collapses onto a line, so that ``everything below $v$'' is one interval
$[\tin(v),\tout(v))$ (\thmref{c24-range}) and ``is $u$ above $v$'' is two comparisons
(\thmref{c24-anc}). Read the same tour a different way and you get the path as a stack, the
depth sequence, and the $2n{-}1$ contour whose shallowest node between two marks is their lowest
common ancestor (\thmref{c24-dist}).

\medskip
Carry three reflexes forward:
\begin{enumerate}[nosep, leftmargin=1.6em]
  \item \emph{Flatten first, then decide the structure.} Whenever a problem asks about a whole
        subtree many times, run one DFS for $\tin/\size$ and reduce it to a range -- the range
        engine (prefix / Fenwick / segment) is a separate, already-solved choice.
  \item \emph{Push and undo.} A root-path aggregate is a value added on entry and subtracted on
        exit; forget the pop and every sibling branch is inflated (\propref{c24-sweep}).
  \item \emph{Distance is LCA in disguise.} $\dist(u,v)=\depth u+\depth v-2\depth{\lca}$; locate
        the LCA by the tour's min-depth window and hand the rest to \chref{c11}.
\end{enumerate}

\medskip
And three habits, drilled until boring: write $G$'s $\tin$, $\tout$, $\size$ and bracket string
by hand before trusting a solution; keep straight the two clocks (both-tick for ancestry,
entry-only for indexing) and know which one your interval endpoints came from; and never answer
a subtree query by re-walking the subtree once you own the flat array.

\medskip
\textbf{Next:} \chref{c25} keeps the DFS-order intuition but flips the arrow -- instead of laying
a given tree onto a line, you \emph{rebuild} the tree from data its walks left behind (two
orders, a bracket string, a parent array), the constructive mirror of everything in this chapter.

\begin{center}\textcolor{softpurple}{\rule{0.4\linewidth}{0.6pt}\quad $\blacksquare$\quad \rule{0.4\linewidth}{0.6pt}}\end{center}
\end{lastword}
